\exercisesection{3}

\begin{enumerate}
\def\labelenumi{\arabic{enumi}.}
\tightlist
\item
  \begin{enumerate}
  \def\labelenumii{(\alph{enumii})}
  \tightlist
  \item
    For the direct sum of a family \((\mathbf{E}_{\lambda})\) of A-modules to be faithfully flat, it is sufficient that each of the \(E_{\lambda}\) be flat and that at least one of them be faithfully flat.
  \end{enumerate}
\end{enumerate}

\begin{enumerate}
\def\labelenumi{(\alph{enumi})}
\setcounter{enumi}{1}
\tightlist
\item
  Deduce from (a) that, if A is a simple ring, every non-empty A-module is faithfully flat. Is the result true for semisimple rings?
\end{enumerate}

\begin{enumerate}
\def\labelenumi{\arabic{enumi}.}
\setcounter{enumi}{1}
\item
  Let \((p_n)\) be the strictly increasing sequence of prime numbers and A the product ring \(\prod_{n} \mathbf{Z} / p_n \mathbf{Z}\) . Show that the direct sum E of the \(\mathbf{Z} / p_n \mathbf{Z}\) is a faithful projective A-module which is not faithfully flat (observe that E is an ideal of A such that \(\mathrm{E}^2 = \mathrm{E}\) ).
\item
  Let A be a right coherent ring (§ 2, Exercise 12). For a product of left A-modules to be faithfully flat, it is sufficient that each of them be flat and that at least one of them be faithfully flat.
\end{enumerate}

Deduce that, if A is a coherent commutative ring, the ring of formal power series \(A[[X_{1},\ldots,X_{n}]]\) is a faithfully flat A-module.

\begin{enumerate}
\def\labelenumi{\arabic{enumi}.}
\setcounter{enumi}{3}
\item
  Let A be a simple algebra over a commutative field K and B a subalgebra of A which is semisimple but not simple. Show that A is a faithfully flat (right or left) B-module, but that there exist right B-modules E which are not faithfully flat, whilst \(E \otimes_{B} A\) is always faithfully flat (Exercise 1).
\item
  Let A be a commutative ring and M a flat A-module containing a submodule N which is not a flat module (cf.~§ 2, Exercise 3). Let B (resp. C) be the A-module A ⊕ N (resp. A ⊕ M) in which multiplication is defined by \((a, x)(a', x') = (aa', ax' + a'x)\) ; then B is not a flat A-module, but the B-module C is a faithfully flat A-module and consequently B satisfies the conditions of Proposition 8 of no. 5.
\item
  Give an example of an integral domain A and a ring B of which A is a subring, such that B is a flat A-module but there exists an A-module E which is neither projective nor finitely generated, for which \(B \otimes_{A} E\) is a finitely generated free B-module.
\item
  If K is a field, the ring K{[}X{]} and the field K(X) are faithfully flat K-modules, but K(X) is not a faithfully flat K{[}X{]}-module.
\item
  Let \(p\) be a prime number and \(\mathbf{A}\) the subring of \(\mathbf{Q}\) consisting of the fractions \(k / p^n\) , where \(k \in \mathbf{Z}\) , \(n \geqslant 0\) . Show that \(\mathbf{A}\) is a flat \(\mathbf{Z}\) -module and that there exists a \(\mathbf{Z}\) -module \(\mathbf{E}\) which is not flat but where \(\mathbf{A} \otimes_{\mathbf{Z}} \mathbf{E}\) is a flat \(\mathbf{A}\) -module.
\item
  Let A be a commutative ring, B an A-algebra, \((\mathbf{C}_{\lambda})_{\lambda \in L}\) , a family of A-algebras and \(B_{\lambda} = C_{\lambda} \otimes_{A} B\) the tensor product algebra of \(C_{\lambda}\) and B for all \(\lambda \in L\) . Let E be a left B-module. Set \(E_{\lambda} = B_{\lambda} \otimes_{B} E = C_{\lambda} \otimes_{A} E\) ; this is a \((\mathbf{B}_{\lambda}, \mathbf{C}_{\lambda})\) -bimodule. Similarly, if F is a right B-module, set
\end{enumerate}

\[
\mathbf {F} _ {\lambda} = \mathbf {F} \otimes_ {\mathbf {B}} \mathbf {B} _ {\lambda} = \mathbf {F} \otimes_ {\mathbf {A}} \mathbf {C} _ {\lambda},
\]

which is a \((\mathbf{C}_{\lambda},\mathbf{B}_{\lambda})\) -bimodule.

\begin{enumerate}
\def\labelenumi{(\alph{enumi})}
\item
  Show that the \((\mathbf{C}_{\lambda},\mathbf{C}_{\lambda})\) -bimodule \(\mathbf{F}_{\lambda}\otimes_{\mathbf{B}_{\lambda}}\mathbf{E}_{\lambda}\) is isomorphic to \((\mathbf{F}\otimes_{\mathbf{B}}\mathbf{E})\otimes_{\mathbf{A}}\mathbf{C}_{\lambda}\) .
\item
  Show that, if E is a flat (resp. faithfully flat) B-module, each of the \(E_{\lambda}\) is a flat (resp. faithfully flat) \(B_{\lambda}\) -module. The converse is true if we assume further that \(\bigoplus C_{\lambda}\) is a faithfully flat A-module.
\item
  Show that, if \(L\) is finite, each of the \(E_{\lambda}\) is a finitely generated projective \(B_{\lambda}\) -module and \(\bigoplus_{\lambda \in L} C_{\lambda}\) is a faithfully flat A-module, then \(E\) is a finitely generated projective B-module (use Proposition 12 of no. 6).
\end{enumerate}

\begin{enumerate}
\def\labelenumi{\arabic{enumi}.}
\setcounter{enumi}{9}
\tightlist
\item
  \begin{enumerate}
  \def\labelenumii{(\alph{enumii})}
  \tightlist
  \item
    Let \(\rho: A \to B\) be a ring homomorphism. Show that, for every left ideal \(a\) of \(A\) which is a left annihilator of a subset \(M\) of \(A\) , \(\rho^{-1}(Ba) = a\) .
  \end{enumerate}
\end{enumerate}

\begin{enumerate}
\def\labelenumi{(\alph{enumi})}
\setcounter{enumi}{1}
\tightlist
\item
  Deduce from (a) an example of a homomorphism \(\rho: A \to B\) such that right A-module \(B\) is not flat but \(\bar{\rho}^{-1}(Ba) = a\) holds for every left ideal \(a\) of \(A\) (cf.~§ 2, Exercise 17 and Algebra, Chapter VIII, § 3, Exercise 11 and § 2, Exercise 6 and Chapter IX, § 2, Exercise 4).
\end{enumerate}

